% ============================================================
% tab/tabla_comparativa.tex
%
% COLUMNAS OBLIGATORIAS. No sustituir, no reordenar, no añadir.
%
% CRITERIO DE INCLUSIÓN: trabajos que resuelven una tarea sobre
% un conjunto de datos y reportan una métrica. Las referencias
% metodológicas se discuten en el texto y no figuran aquí.
%
% Se divide en dos partes por volumen. La segunda retrocede el
% contador para conservar un único número de tabla.
%
% REGLA DE tabularx: la suma de los pesos W debe ser igual al
% número de columnas W. Aquí son 5 y suman 5.00.
% ============================================================

\begin{table*}[!tp]
\centering
\caption{Resumen comparativo de trabajos relacionados. Los valores están
transcritos de la fuente. Los conjuntos difieren en tamaño, composición,
juego y versión, por lo que la comparación directa requiere cautela.}
\label{tab:comparativa}
\footnotesize
\setlength{\tabcolsep}{3.5pt}
\begin{tabularx}{\textwidth}{@{}
    P{1.9cm}
    >{\centering\arraybackslash}p{0.6cm}
    W{0.95}
    W{1.00}
    W{0.70}
    W{0.95}
    W{1.40}
  @{}}
\toprule
\textbf{Referencia} & \textbf{Año} & \textbf{Técnica} &
\textbf{Conjunto de datos / entorno} & \textbf{Métrica principal} &
\textbf{Resultado reportado} & \textbf{Limitación identificada}\\
\midrule

Janusz et al. \cite{janusz_hearthstone_2017} & 2017 &
Competencia abierta; la referencia de los organizadores es una red de dos
capas ocultas sobre representación tabular &
\emph{Hearthstone}: estados intra-partida; conjuntos de entrenamiento y
prueba separados por partida y por mazo &
AUC &
0.7846 la referencia; 0.8019 la mejor solución. Al transferir a estados
generados por agentes de búsqueda, el AUC cae de $\approx$0.79 a
$\approx$0.75 &
Estados generados por agentes que eligen acciones al azar, población no
representativa; no se reporta calibración; una fuga involuntaria a mitad de
competencia obligó a reasignar parte del conjunto de prueba
\\
\addlinespace[3pt]

Hodge et al. \cite{hodge_winprediction_2021} & 2021 &
Regresión logística, bosques aleatorios y potenciación del gradiente sobre
ventanas deslizantes del estado &
DotA 2: repeticiones mixtas y profesionales &
Exactitud &
Exactitudes en el rango aproximado de 0.70 a 0.80, evaluando en el punto
medio de la partida &
Los conjuntos de la literatura no son comparables entre sí; el flujo en vivo
provee un subconjunto de las características presentes en las repeticiones
\\
\addlinespace[3pt]

Guo et al. \cite{guo_temporalshift_2021} & 2021 &
Revisión sistemática de estrategias de mitigación del cambio temporal &
15 estudios de predicción clínica retenidos de 4\,457 referencias
examinadas; regresión logística en 13 de 15 &
Cambio en calibración y en discriminación &
Deterioro de calibración en 11 de los 11 estudios que la midieron;
deterioro de discriminación en solo 3 de 12 &
Dominio clínico, citado por la asimetría que mide y por la terminología que
estandariza; sin metaanálisis; cobertura incompleta de actas de congresos
\\
\addlinespace[3pt]

Tian et al. \cite{tian_hero_2022} & 2022 &
Red recurrente con características de alineación y de combate &
Honor of Kings: 27\,256 partidas, 496\,342 marcos &
Exactitud &
81.8\,\% en el nodo de 7 min según la tabla detallada; el resumen del
artículo reporta 84.7\,\% &
Datos de un solo día y jugadores de nivel elevado; partición descrita a
nivel de marcos sin agrupación explícita por partida; no evalúa calibración
\\
\addlinespace[3pt]

Miernik y Kowalski \cite{miernik_evolving_2022} & 2022 &
Evolución de funciones de evaluación lineales y basadas en árboles &
Legends of Code and Magic &
Tasa de victoria &
60.2\,\% promedio global para la variante de mejor desempeño &
No estima probabilidades de forma supervisada; posible especialización
respecto del oponente empleado para evaluar
\\
\addlinespace[3pt]

Xenopoulos et al. \cite{xenopoulos_winprob_2022} & 2022 &
Máquinas de potenciación explicables, un modelo aditivo generalizado con
atribuciones nativas &
CSGO: tres poblaciones de distinto nivel de habilidad &
Pérdida logarítmica y error de calibración esperado &
Modelos bien calibrados dentro de su propio entorno. El entrenado sobre la
población de menor nivel presenta mayor pérdida logarítmica y peor
calibración al evaluarse fuera &
Un solo juego; el modelo interpretable podría ser superado por otros
métodos; no es posible identificar a los jugadores que incurren en trampa
\\
\addlinespace[3pt]

Bia\l ecki et al. \cite{bialecki_sc2egset_2023} & 2023 &
Regresión logística, máquina de vectores de soporte y potenciación del
gradiente con hiperparámetros próximos a los predeterminados &
SC2EGSet: repeticiones de torneos de StarCraft II; 11\,115 partidas
reducidas a 22\,230 muestras promediadas; 17 características económicas &
Exactitud &
0.9071 con dos jugadores; 0.8924 y 0.8916 los otros dos modelos; 0.8488 con
un solo jugador &
Solo jugadores de nivel de torneo, sin análisis entre niveles de habilidad;
los autores describen sus experimentos como ilustrativos; no se documenta
agrupación por partida en los pliegues
\\

\bottomrule
\end{tabularx}
\end{table*}


% ============================================================
% PARTE 2 DE 2 — el contador retrocede para mantener un número.
% ============================================================
\begin{table*}[!tp]
\addtocounter{table}{-1}
\centering
\caption{Resumen comparativo de trabajos relacionados (continuación).}
\label{tab:comparativa-cont}
\footnotesize
\setlength{\tabcolsep}{3.5pt}
\begin{tabularx}{\textwidth}{@{}
    P{1.9cm}
    >{\centering\arraybackslash}p{0.6cm}
    W{0.95}
    W{1.00}
    W{0.70}
    W{0.95}
    W{1.40}
  @{}}
\toprule
\textbf{Referencia} & \textbf{Año} & \textbf{Técnica} &
\textbf{Conjunto de datos / entorno} & \textbf{Métrica principal} &
\textbf{Resultado reportado} & \textbf{Limitación identificada}\\
\midrule

Chitayat et al. \cite{chitayat_meta_2023} & 2023 &
Representación de personajes por atributos de diseño, agrupada con K-Means,
sobre redes neuronales &
DotA 2: versiones 7.27 a 7.33, 61\,254 partidas &
AUC &
0.85 en versiones conocidas; 0.85 y 0.86 en dos versiones posteriores no
empleadas para entrenar &
Solo partidas profesionales; predice conteos finales a partir de duración y
composición, no probabilidad desde un estado intermedio
\\
\addlinespace[3pt]

Bertram et al. \cite{bertram_cardrepresentations_2024} & 2024 &
Redes siamesas sobre representaciones generalizadas de carta construidas
con atributos, metadatos e imágenes &
Selección de cartas de Magic: The Gathering; un conjunto completo de cartas
reservado para evaluación &
Acierto \emph{top-1} &
67.20\,\% sobre conjuntos vistos frente a 55.44\,\% sobre el conjunto no
visto; la codificación por identidad queda indefinida sobre cartas nuevas &
La tarea es predecir una elección humana y no un desenlace de partida, por
lo que solo es análoga en cuanto a generalización de contenido; las
elecciones humanas no son consistentes, lo que acota el acierto alcanzable
\\
\addlinespace[3pt]

Saravanan y Guzdial \cite{saravanan_metadiscovery_2024} & 2024 &
Agente de refuerzo y agente heurístico, con constructor de equipos &
Pok\'emon Showdown, el videojuego competitivo; cuatro escenarios de
prohibición de personajes &
Solapamiento, distancia de edición y correlación de rangos &
Solapamiento elevado entre el conjunto de estrategias descubierto y el real,
pero correlación de rangos débil y no significativa &
La correlación de rangos no significativa indica que el conjunto descubierto
difiere del real pese al solapamiento; es el videojuego y no el juego de
cartas
\\
\addlinespace[3pt]

Kamal et al. \cite{kamal_moba_2025} & 2025 &
Revisión sistemática con metodología PRISMA &
35 estudios sobre deportes electrónicos de arena multijugador &
Métricas heterogéneas &
Sin resultado único; se reportan, entre otros, AUC de 0.97 y exactitud
cercana al 93\,\% &
Alta heterogeneidad entre tareas, conjuntos y protocolos; generalización
limitada entre regiones, periodos y poblaciones
\\
\addlinespace[3pt]

Karbalaie et al. \cite{karbalaie_participantaware_2026} & 2026 &
Diez clasificadores bajo cuatro diseños de validación, incluida la
estratificada y la agrupada por participante &
623 ensayos de 72 participantes; 624 características reducidas a un
conjunto fijo de 10 &
Exactitud y brecha entrenamiento--prueba &
0.91 bajo validación estratificada frente a 0.66 dejando un participante
fuera; los modelos lineales difieren en 0.03 o menos entre ambos protocolos &
Dominio de biomecánica, citado por su argumento estadístico y no como
resultado del dominio de juegos; una única tarea binaria; calibración
explícitamente fuera de alcance
\\
\addlinespace[3pt]

Angliss et al. \cite{angliss_vgcbench_2026} & 2026 &
Clonación de comportamiento y aprendizaje por refuerzo multiagente &
Pok\'emon VGC, el videojuego competitivo; 72 composiciones reservadas fuera
de distribución &
Tasa de victoria y evaluación cruzada &
0.669 del agente entrenado con 64 composiciones frente al entrenado con una,
sobre las composiciones no vistas &
El desempeño en configuraciones familiares disminuye al crecer la
diversidad; no es una tarea de estimación de probabilidad
\\

\bottomrule
\end{tabularx}
\begin{minipage}{\textwidth}
\vspace{3pt}
\footnotesize
\textit{Nota.} Las referencias de carácter metodológico incluidas en la
revisión ---formalización de la fuga de datos, catálogos de escenarios de
fuga y revisiones de calibración y de modelos tabulares--- no figuran en
esta tabla, dado que no resuelven una tarea de predicción sobre un conjunto
de datos comparable. Sus aportes se discuten en el texto.
\end{minipage}
\end{table*}